\textbf{COLUNI 2020 -- Questão 06}\par\medskip Leia o poema.

\par

Enredo para um tema

\par

Ele me amava, mas não tinha dote,

só os cabelos pretíssimos e uma beleza

de príncipe de estórias encantadas.

Não tem importância, falou a meu pai,

se é só por isto, espere.

Foi-se com uma bandeira

e ajuntou ouro pra me comprar três vezes.

Na volta me achou casada com D. Cristóvão.

Estimo que sejam felizes, disse.

O melhor do amor é sua memória, disse meu pai.

Demoraste tanto, que... disse D. Cristóvão.

Só eu não disse nada,

nem antes, nem depois.

\par

PRADO, A. \textit{Bagagem}. 29ª ed. Rio de Janeiro. Record. 2019. p. 91.

\par

De acordo com o poema acima, assinale a afirmativa CORRETA:
\begin{enumerate}
\item Em: “\underline{Ele me amava, mas não tinha dote},”, estabelece-se, no período, a ideia de explicação.
\item Em: “[...] cabelos \underline{pretíssimos}”, o superlativo enfatiza uma característica dos cabelos.
\item Em: “\underline{Foi-se} com uma bandeira”, o se tem o mesmo sentido que em: “\underline{se} é só por isto, espere”.
\item Em: “[...] se é só por \underline{isto}, espere.”, o pronome isto refere-se a príncipe de estórias encantadas.
\end{enumerate}
